\documentclass[a4paper]{article}     
\usepackage{amsfonts}
\usepackage{amsmath}
\usepackage{amssymb}
\usepackage{graphicx}

\renewcommand{\arraystretch}{2.5}

\begin{document}

\noindent \large Auteur: Elias Alsa Ati                         \\
\noindent \large Studierichting: Informatica                    \\
\noindent \large Oefeningenreeks 1D nr. 6.9 \small\cite{a}      \\

\medskip

\normalsize

$$\left( \left(2 + i\right)z^3 - iz^2 + 3\right) \div \left(\left(1 - i\right)z^2 - z + 1\right)$$

$\begin{array}{rccccc|c}
&   &\left(2 + i\right)z^3  &    - iz^2                     &   + 0z                                    &       + 3             &   \left(1 - i\right)z^2 - z + 1                       \\  \cline{1-7}
& - &\left(2 + i\right)z^3  &    - \dfrac{1 + 3i}{2} z^2    &   + \dfrac{1 + 3i}{2} z                   &        |              &    \dfrac{1+3i}{2}z + \dfrac{1}{2}i \rightarrow Q     \\  \cline{3-5}
&   &          0            &    +  \dfrac{1+i}{2}z^2       &   - \dfrac{1 + 3i}{2}z                    &       + 3             &                                                       \\
& - &                       &       \dfrac{1+i}{2}z^2       &   - \dfrac{1}{2}iz                        &       +\dfrac{1}{2}i  &                                                       \\  \cline{4-6} 
&   &                       &          R \leftarrow         &    \left(-\dfrac{1}{2} -i\right)z         &   + 3 - \dfrac{1}{2}i &
\end{array}$ \\

\textbf{Resultaat:}\\

\begin{align*}
    A(z) &= B(z) \cdot Q(z) + R(z) \\
    A(z) &= \left(\left(1 - i\right)z^2 - z + 1\right) \cdot \left(\dfrac{1+3i}{2}z + \dfrac{1}{2}i\right) + \left(\left(-\dfrac{1}{2}-i\right)z + 3 -\dfrac{1}{2}i\right)
\end{align*}

\textbf{KladBlok:}  \\

Berekening van het eerste quotiëntterm:\\

\begin{align*}
    \dfrac{\left(2 + i\right)z^3}{\left(1 - i\right)z^2}    &= \dfrac{\left(2 + i\right)}{\left(1 - i\right)} z                 \\
                                                            &= \left(2 + i\right) \cdot \dfrac{\left(1 + i\right)}{1^2 + 1^2}z  \\
                                                            &= \dfrac{2 + 2i + i + i^2}{2}z                                     \\
                                                            &= \dfrac{1 + 3i}{2}z
\end{align*}

Berekening van het 2de quotiëntterm:\\

\begin{align*}
    \dfrac{\dfrac{1+i}{2}z^2}{\left(1 - i\right)z^2}        &= \dfrac{1 + i}{2\left(1 - i\right)}                                       \\
                                                            &= \dfrac{\left(1 + i\right)\left(1 + i\right)}{2\left(1^2 + 1^2\right)}    \\
                                                            &= \dfrac{1 + 2i + i^2}{4}                                                  \\
                                                            &= \dfrac{2}{4}i = \dfrac{1}{2}i
\end{align*}

\begin{thebibliography}{999}
\bibitem{a} 
\end{thebibliography}
\end{document} 